\Needspace{42\baselineskip}
\section{Grupoide de la extensión bidireccional y álgebra operatorial}
\label{sec:cpe-grupoide-two-way}

Sea \(X_F^{\leftrightarrow}\) la extensión natural del estado enriquecido y
\(\sigma\) su homeomorfismo de desplazamiento. Se define el grupoide de
transformación
\begin{equation}
 \mathcal G_{\rm TPK}
 :=X_F^{\leftrightarrow}\rtimes_\sigma\mathbb Z,
 \label{eq:cpe-grupoide-tpk}
\end{equation}
con
\[
 s(x,n)=x,
 \quad r(x,n)=\sigma^n x,
 \quad (\sigma^nx,m)(x,n)=(x,m+n),
 \quad (x,n)^{-1}=(\sigma^nx,-n).
\]

\begin{theorem}[Completación grupoidal del transporte bidireccional]
\label{thm:cpe-grupoide-tpk}
El grupoide \(\mathcal G_{\rm TPK}\) es localmente compacto, Hausdorff,
segundo numerable, étale y amenable. Para la medida proyectiva invariante y
no atómica \(\mu_{\rm HMT}\), si el desplazamiento es ergódico y
esencialmente libre,
\[
 \mathcal M_{\rm TPK}
 :=L^\infty(X_F^{\leftrightarrow},\mu_{\rm HMT})
 \rtimes_\sigma\mathbb Z
\]
es el factor hiperinfinito de tipo \(II_1\).
\end{theorem}

\begin{proof}
La acción de un grupo discreto numerable por homeomorfismos sobre un compacto
metrizable produce un grupoide étale, localmente compacto, Hausdorff y
segundo numerable. La amenabilidad de \(\mathbb Z\) implica la amenabilidad
del grupoide. La medida proyectiva es invariante por el desplazamiento.
Ergodicidad y libertad esencial hacen factorial el producto cruzado; la no
atomicidad excluye el caso matricial de tipo I. La traza inducida por
\(\mu_{\rm HMT}\) es finita y fiel, y la amenabilidad produce
hiperfinitud. Por tanto, el factor es de tipo \(II_1\).
\end{proof}

La clasificación modular de tipo \(III\) pertenece a una medida
cuasi-invariante distinta, con cociclo de Radon--Nikodym declarado. No se
deduce de la mera presencia de una dinámica KMS y no reabre la completación
grupoidal anterior.
